\section{Related Work}

\textbf{CVSS, EPSS, and exploit prediction.} CVSS alone poorly predicts exploitation
\cite{allodi2014comparing}, motivating data-driven scoring including the EPSS lineage
\cite{jacobs2021epss} and earlier signal-based approaches \cite{sabottke2015vulnerability}.
We consume these as features/baselines rather than proposing a new predictor.

\textbf{Risk-based vulnerability management and commercial RBVM.} Commercial RBVM
products --- Microsoft Defender Vulnerability Management [VERIFY: Microsoft Defender
VM docs], Tenable VPR/ACR \textbf{[VERIFY: Tenable VPR/ACR docs]}, Qualys TruRisk [VERIFY:
Qualys TruRisk docs], and Cisco/Kenna \textbf{[VERIFY: Cisco/Kenna docs]} --- combine exploit
signals with proprietary context. They are acknowledged as industry practice but
are \textbf{not} benchmarked here: their models are closed and not independently
reproducible.

\textbf{Machine-learning vulnerability prioritization.} Research systems such as Deep
VULMAN \textbf{[VERIFY: Deep VULMAN source]}, VulRG \textbf{[VERIFY: VulRG source]}, VulnScore
\textbf{[VERIFY: VulnScore source]}, and V-REx \textbf{[VERIFY: V-REx source]} apply learning and
risk/graph modeling to prioritization. They advance modeling quality but typically
report a single tuned configuration on non-reproducible data and do not emphasize
per-decision audit evidence or an open freeze/verify protocol. We position this
work as complementary infrastructure and do \textbf{not} claim to be first.

\textbf{Resource-constrained remediation.} Our scheduler models bounded per-window
capacity, blackout windows, approval policies, and risk acceptance, so strategies
are evaluated jointly with the operational constraints that shape real outcomes.

\textbf{Public-sector compliance and auditability.} Guidance such as NIST SP 800-40
\cite{nist80040r4} and SP 800-53 \cite{nist80053r5} and CISA directives \cite{cisa_bod2201}
frames remediation as a documented, reviewable process; our append-only,
hash-chained audit log targets the per-decision evidence such processes call for.
